\documentclass[10pt,a4paper]{amsart}
\usepackage[margin=19mm]{geometry}
\usepackage{amsmath,amssymb,mathtools}
\usepackage[scr=rsfs,cal=euler]{mathalfa}
\usepackage{xfrac}
\usepackage[unicode,hidelinks,bookmarks=false]{hyperref}
\newcommand{\ie}{i.e.\ }
\newcommand{\cf}{cf.\ }
\newcommand{\resp}{respectively\ }
\newcommand{\Subsection}{\subsection}
\providecommand{\nobreakdash}{\nobreak\hbox{-}}
\setlength{\emergencystretch}{2em}
\multlinegap0pt
\usepackage{amssymb}
\usepackage{xcolor}
\usepackage{tikz}
\newtheorem{proposition}{Proposition}
\newtheorem{corollary}{Corollary}
\newtheorem{lemma}{Lemma}
\newtheorem{claim}{Claim}
\newcommand{\br}{{\mathscr{B}}}
\newcommand{\per}{\mathcal{H}}
\title{Relative accessibility for graphs}
\author{Joseph Paul MacManus}
\date{Source: arXiv:2605.12629v1 (CC BY 4.0)}
\begin{document}
\maketitle

\noindent\emph{Source and licence.} Relative accessibility for graphs. Version arXiv:2605.12629v1 (CC BY 4.0), distributed by arXiv under the Creative Commons Attribution 4.0 International licence, http://creativecommons.org/licenses/by/4.0/, as recorded in the arXiv licence metadata for this exact version and shown on the abstract page https://arxiv.org/abs/2605.12629v1. Full licence notice: \texttt{licenses/arxiv-2605.12629-CC-BY-4.0.txt}.\par\smallskip

\noindent\textit{Editorial scope.} Counted unit: the article's lemma claim:I-fg-implies-H-fg (source0.tex lines 1458--1460). The article's complete original proof of it is delivered with it (source0.tex lines 1462--1524, one \texttt{proof} environment of 63 lines). No mathematics is written, repaired or re-derived: the statement and the proof are the article's own lines, in the article's own notation, and no retained line is substituted or re-flowed.  Retained uncounted as the source-local context the proof needs: lines 424--430; lines 1194--1196; lines 1317--1319; lines 1337--1347.  Nothing was omitted from any retained span, so the package carries no omission record.  No retained line contains a citation, so the wrapper carries no bibliography and no bibliography entry.  No retained line contains a figure environment, an image-inclusion command or drawing code, so no image is delivered and the package declares no asset dependency.  Editorial formatting only, in addition: an AMS \texttt{amsart} preamble that restates the article's own declarations and macros used by the retained lines, namely the environments \texttt{proposition}, \texttt{corollary}, \texttt{lemma}, \texttt{claim} on separate counters, together with the standard packages those lines need.  The retained environments print in this document as Proposition 1, Corollary 1, Lemma 1, Claim 1 in order of appearance, whereas the article prints the same items with the numbering of its own full document.  The complete native source of the article as distributed in the arXiv e-print for this exact version is bundled unchanged as \texttt{sources/200431/source0.tex}.
\begin{proposition}\label{prop:g-finite-criterion}
    Let $\Gamma$ be a connected, locally finite graph. Let $M \subset \br(\Gamma)$ be a subring. Let $S \subset M$ be such that
    $$
    \text{$\exists b \in M$ separating $\omega_1$, $\omega_2$ $\iff$ $\exists b' \in S$ separating $\omega_1$, $\omega_2$ }.
    $$
    Then for every $b \in M$ there exists $b' \in \langle S \rangle$, the subring generated by $S$, such that $b+b'$ is finite.
\end{proposition}

\begin{corollary}\label{cor:gen-set-upgrade-rel-cuts}
    Let $\Gamma$ be a connected $G$-graph and $\per$ a tame, $G$-invariant peripheral system. Suppose $\br_\per(\Gamma)$ is $G$-finitely generated as a ring. Then $\br_\per(\Gamma)$ is a finitely generated $G$-module. 
\end{corollary}

    \begin{lemma}\label{claim:finding-the-special-cut}
        Let $H \in \mathcal T$. Let $U \supset \Lambda(H)$ be an open neighbourhood of the limit set of $H$ in $\Omega(\Gamma)$. Then there exists $b \in \br_\per(\Gamma)$ such that for all $\omega \in \Omega(\Gamma) \setminus U$, we have that $b$ separates $\omega$ from $\Lambda(H)$. 
    \end{lemma}

    \begin{lemma}\label{claim:technical-boundary-claim}
        Let $H_1, \ldots, H_k \in \mathcal T$, with $\{\lambda_{j,1}, \lambda_{j,2}\} = \Lambda(H_j)$. Let $W_1 \sqcup W_2 = \Omega(\Gamma)$ be a bipartition into clopen sets, such that $\lambda_{j,i} \in W_i$ for all $j$, $i$. Then there exists pairwise disjoint, open neighbourhoods $U_{j,i} \ni \lambda_{j,i}$ in $\Omega(\Gamma)$, with $U_{j,i} \subset W_i$, such that the following holds:
            
            For all $i \in \{1,2\}$, $j \in \{1,\ldots, k\}$, $p \in U_{j,i} \setminus \Lambda(H_j)$, there exists $g \in G_{H_j}$ such that:
            \begin{enumerate}
                \item\label{itm:dynamics-property-1} $gp \in W_i \setminus \bigcup _\ell U_{\ell, i}$, and

                \item\label{itm:dynamics-property-2} $gU_{j,i+1} \cap W_i= \emptyset$,
            \end{enumerate}
            where addition within the indices is taken modulo 2.
    \end{lemma}

    \begin{lemma}\label{claim:I-fg-implies-H-fg}
        If $\br_{\mathcal I}(\Gamma)$ is a finitely generated $G$-module, then so is $\br_\per(\Gamma)$. 
    \end{lemma}

    \begin{proof}
        Let $S \subset \br_{\mathcal I}(\Gamma)$ be a finite subset so that $G \cdot S$ additively generates $\br_{\mathcal I}(\Gamma)$.  Bipartition $S$ into two sets $S = S_{\mathrm{ell}} \sqcup S _{\mathrm{hyp}}$, where
        $$
        S_{\mathrm{ell}} := \{b \in S : \text{$b$ is $\per$-elliptic}\}, \ \ \ S_{\mathrm{hyp}} := \{b \in S : \text{$b$ is not $\per$-elliptic}\}.
        $$
        Note that if $S_{\mathrm{hyp}} = \emptyset$, then there is nothing to prove, so let's assume this set is non-empty. For each $b \in S_{\mathrm{hyp}}$, let $H_1^{b} \ldots, H_{k_b}^b \in \mathcal T$ denote the (necessarily finite and non-empty) collection of $\mathcal T$-peripherals which have infinite intersection with both $b$ and $b^\ast$. 

        Every $b \in S_{\mathrm{hyp}}$ induces a non-trivial bipartition of $\Omega(\Gamma)$ into clopen sets $W_1^b \sqcup W_2^b$. Let $\Lambda(H_j^b) = \{\lambda_{j,1}^b, \lambda_{j,2}^b\}$, indexed so that $\lambda^b_{j,i} \in W_i^b$. 
        We now apply Claim~\ref{claim:technical-boundary-claim}, and for each $b \in S_{\mathrm{hyp}}$ we obtain pairwise disjoint neighbourhoods $U_{j,i}^b$ of $\lambda_{j,i}^b$, such that $U_{j,i}^b \subset W_i^b$, and such that for all $j \in \{1,\ldots, k_b\}$, and all $p \in U_{i,j}^b$, there exists $g \in G_{H_j^b}$ such that:
            \begin{enumerate}
                \item $gp \in W_1^b \setminus \bigcup _\ell U_{\ell,1}^b$, and

                \item $gU_{j,2}^b \cap W_1^b = \emptyset$.
            \end{enumerate}
        Now, apply Claim~\ref{claim:finding-the-special-cut}, and for each $j \in \{1, \ldots, k_b\}$ find $c^b_j \in \br_\per(\Gamma)$ such that for all $\omega \in \Omega(\Gamma) \setminus (U_{j,1}^b \cup U_{j,2}^b)$, we have that $c_j^b$ separates $\omega$ from $\Lambda(H_j^b)$. These are oriented in such a way that $H_j^b \cap c_j^b$ is finite. 

        Now, consider the set 
        $$
        R = S_{\mathrm{ell}} \cup \Big\{c_j^b : b \in S_{\mathrm{hyp}}, \ j = 1, \ldots, k_b  \Big\} \cup \Bigg\{ b \cap \bigcap_{j=1}^{k_b} c_j^b : b \in S_{\mathrm{hyp}}  \Bigg\}. 
        $$
        Note that $R$ is finite. We have the following two claims.

        \begin{claim}
            We have that $R \subset \br_\per(\Gamma)$.
        \end{claim}

        \begin{proof}
            Indeed, it is sufficient to show that $a = b \cap \bigcap_{j=1}^{k_b} c_j^b $ is $\mathcal T$-elliptic for every $b \in S_{\mathrm{hyp}}$. Suppose some $H \in \mathcal T$ has infinite intersection with both $a$ and $a^\ast$. Then $H$ has infinite intersection with $b$ and with every $c_j^b$. Each $c_j^b$ is $\mathcal H$-elliptic, so since $H \cap a^\ast$ is infinite, we must have that $H \cap b^\ast$ is infinite. But then $H = H_j^b$ for some fixed $j$. But this contradicts the fact that $H \cap c_j^b$ is infinite, which is false by construction. It follows that $R \subset \br_\per(\Gamma)$. 
        \end{proof}

        \begin{claim}
             For all distinct pairs  $\omega_1, \omega_2 \in \Omega(\Gamma)$, either $\omega_1, \omega_2 \in \Lambda(H)$ for some $H \in \per$, or there exists $a \in G \cdot R$, such that $a$ separates $\omega_1$ and $\omega_2$.
        \end{claim}

        \begin{proof}
            
        
        
        Fix $\omega_1$ and $\omega_2$, and assume without loss of generality that they do not lie in the limit set of some common peripheral. After translating, we may assume that $\omega_1$ and $\omega_2$ are separated by some fixed $b \in S$. Assume without loss of generality that $b \in S_{\mathrm{hyp}}$, and that $\omega_i \in W_i^b$. We have a few cases to consider. 

        \textbf{Case 1.} It could be the case that some $c_j^b$ separates $\omega_1$ and $\omega_2$, in which case we are done. 

        \textbf{Case 2.} Suppose that no $c_j^b$ separates $\omega_1$ and $\omega_2$, and that both $\omega_1$ and $\omega_2$ are contained in (the limit set of) the intersection
        $
        \bigcap_{j=1}^{k_b} c_j^b
        $.
        Then it is easy to see that 
        $
        b \cap \bigcap_{j=1}^{k_b} c_j^b 
        $
        separates $\omega_1$ from $\omega_2$. In this case, we are done. 

        \textbf{Case 3.} Suppose that no $c_j^b$ separates $\omega_1$ and $\omega_2$, but the hypotheses of Case 2 do not hold. Note that, by construction, the limit sets of the $(c_{j}^b)^\ast$ are pairwise disjoint, since
        $$
        \Lambda(c^b_j) \cap (U_{j,1}^b \cup U_{j,2}^b) = \emptyset.
        $$
        In particular, in this case we must have that $\omega_i \in U_{j,i}^b$ for some fixed $j \in \{1, \ldots, k_b\}$. By assumption, we do not have that both $\omega_i$ lie in $\Lambda(H_j^b)$, so let's say without loss of generality that $\omega_1$ does not lie in $\Lambda(H_j^b)$. Now, applying Claim~\ref{claim:technical-boundary-claim}, we have by construction that there exists $h \in G_{H_j^b}$ such that $h \omega_1 \in W_1^b \setminus \bigcup _\ell U_{\ell,1}^b$, and $h \omega_2 \in W_2$. Note that $h\omega_1$, $h\omega_2$ are still separated by $b$. By replacing $\omega_1$ and $\omega_2$ with their $h$-translates, we now have that we must be in either Case 1 or Case 2. In particular, $\omega_1$ and $\omega_2$ are separated by some $G$-translate of some element in $R$. 

        These three cases are exhaustive, and so the claim follows.
        \end{proof}
        
        It now follows from Proposition~\ref{prop:g-finite-criterion} and Corollary~\ref{cor:gen-set-upgrade-rel-cuts} that $\br_\per(\Gamma)$ is a finitely generated $G$-module. This concludes the proof of the lemma. 
    \end{proof}

\end{document}
